% 分类目录：dl
\documentclass[12pt, a4paper]{article}
\usepackage[margin=2.5cm]{geometry}
\usepackage{ctex}
\usepackage{amsmath, amssymb}
\usepackage{listings}
\usepackage{xcolor}
\usepackage{tabularx}
\usepackage{hyperref}

\lstset{
    language=Python,
    basicstyle=\ttfamily\small,
    keywordstyle=\color{blue},
    commentstyle=\color{green!50!black},
    stringstyle=\color{red},
    showstringspaces=false,
    breaklines=true,
    numbers=left,
    numberstyle=\tiny\color{gray},
    frame=single
}

\title{知识点精讲：深度学习归一化方法 (Normalization Methods) 详解}
\author{}
\date{}

\begin{document}
\maketitle

\section{归一化的数学统一框架}
对于输入向量 $x$，归一化操作通常遵循以下线性映射：
$$ y = \gamma \frac{x - \mu}{\sqrt{\sigma^2 + \epsilon}} + \beta $$
其中 $\mu$ 和 $\sigma^2$ 是在特定维度内计算的均值和方差，$\gamma$ 和 $\beta$ 为可学习的增益 (Scale) 和偏置 (Shift)。

\section{为什么归一化有效？ (Know the Why)}
\begin{itemize}
    \item \textbf{ICS 假说 vs. 损失曲面平滑化}：BN 最初旨在缓解“内部协变量偏移”(ICS)，但现代研究指出，其核心贡献在于显著平滑了损失曲面的曲率，使得更大的学习率成为可能。
    \item \textbf{权重伸缩不变性}：对于任意标量 $a$，$\text{BN}(aW) = \text{BN}(W)$。这降低了训练对初始化参数幅值的敏感度。
\end{itemize}

\section{主流归一化方法的维度映射}
归一化的本质是在 4D 张量 $(N, C, H, W)$ 中选择不同的计算边界：
\begin{itemize}
    \item \textbf{Batch Norm (BN)}：在 $N$ 维度上归一化。计算各样本相同通道的统计量。
    \item \textbf{Layer Norm (LN)}：在 $C, H, W$ 维度上归一化。计算单个样本内部所有通道的统计量。
    \item \textbf{Instance Norm (IN)}：在 $H, W$ 维度上归一化。常用于风格迁移。
    \item \textbf{Group Norm (GN)}：将 $C$ 分为若干组，在每组内部进行归一化。
\end{itemize}

\section{大模型 (LLM) 时代的 SOTA}

\subsection{RMSNorm (Root Mean Square Normalization)}
这是 Llama 3, Qwen 2 等当代顶级大模型的标准选型。
$$ \text{RMSNorm}(x) = \frac{x}{\sqrt{\frac{1}{n}\sum_{i=1}^n x_i^2 + \epsilon}} \cdot g $$
\textbf{相比 LN 的优势}：移除了均值计算的过程，仅保留根均方缩放。这不仅减少了计算开销，还通过对特征幅值的约束，在深层神经网络中提供了比 LN 更强的数值稳定性。

\subsection{Pre-Norm vs. Post-Norm 架构争论}
\begin{itemize}
    \item \textbf{Post-Norm}：$\text{Norm}(x + F(x))$。虽然能获得更好的泛化上限，但初始化不当极易导致梯度消失/爆炸，难以训练。
    \item \textbf{Pre-Norm}：$x + F(\text{Norm}(x))$。目前的工业界共识，它将残差路径保持为一条“直通的高速公路”，极大提升了极深模型（如 LLM）的训练稳定性。
\end{itemize}

\section{工程细节：SyncBN 与 BatchNorm 陷阱}
\begin{itemize}
    \item \textbf{SyncBN (同步归一化)}：在多卡分布式训练中，BN 若仅在单卡独立计算，会导致等效 Batch Size 过小，统计量偏差大。SyncBN 在多卡间同步统计量，是目标检测等任务的必备。
    \item \textbf{训练 vs 推理}：BN 在推理时必须使用全局滑动平均 (Running Average) 的统计量，而非当前测试样本的均值。
\end{itemize}

\section{全维度对比概览表}
\begin{table}[h]
\centering
\small
\begin{tabularx}{\textwidth}{|l|c|c|X|}
\hline
\textbf{方法} & \textbf{对 Batch Size 依赖} & \textbf{核心维度} & \textbf{典型应用场景} \\
\hline
BN & 极强 (且需大 bs) & $N$ & CNN (ResNet, EfficientNet) \\
\hline
LN & 无 & $C, (H, W)$ & Transformer, RNN (NLP 主流) \\
\hline
RMSNorm & 无 & $C, (H, W)$ & \textbf{Llama 3, GPT-4 (LLM 标准)} \\
\hline
GN & 无 & 组内 $C$ & 小 batch 视觉任务 \\
\hline
\end{tabularx}
\end{table}

\section{PyTorch 工程实现 (Custom RMSNorm)}
\begin{lstlisting}
import torch
import torch.nn as nn

class RMSNorm(nn.Module):
    def __init__(self, dim, eps=1e-6):
        super().__init__()
        self.eps = eps
        self.weight = nn.Parameter(torch.ones(dim))

    def _norm(self, x):
        return x * torch.rsqrt(x.pow(2).mean(-1, keepdim=True) + self.eps)

    def forward(self, x):
        # x.shape: (Batch, SeqLen, HiddenDim)
        return self._norm(x.float()).type_as(x) * self.weight
\end{lstlisting}

\section{参考文献}
\begin{enumerate}
    \item Ioffe, S. (2015). \textit{Batch Normalization}. (Original Paper)
    \item Santurkar, S., et al. (2018). \textit{How Does Batch Normalization Help Optimization?}.
    \item Zhang, B., \& Sennrich, R. (2019). \textit{Root Mean Square Layer Normalization}. (RMSNorm)
\end{enumerate}

\end{document}
